% Section 2: The bridge between forecast loss and decision loss
\subsection{The bridge between forecast loss and decision loss}
\label{sec:bridge}

Consider the newsvendor of Section~\ref{sec:forecasting_rl_intro} in its
simplest binary form. Each day the retailer chooses whether to stock an
extra unit ($a = 1$, Yes) or not ($a = 0$, No). The day is Bad ($y = 1$, a
stockout) with probability $\pi$ or Good ($y = 0$) with probability
$1 - \pi$. Realised utilities are
\begin{equation}
\begin{array}{c|cc}
                  & y = 1 \text{ (Bad)} & y = 0 \text{ (Good)} \\ \hline
a = 1 \text{ (Yes)} & U_{yb} = 0           & U_{yg} = -1 \\
a = 0 \text{ (No)}  & U_{nb} = -4          & U_{ng} = 0
\end{array}
\label{eq:gp_payoff_table}
\end{equation}
The retailer with forecast $\hat\pi$ acts Yes whenever
$\hat\pi\, U_{yb} + (1 - \hat\pi) U_{yg} > \hat\pi\, U_{nb} +
(1 - \hat\pi) U_{ng}$, which rearranges to the threshold rule
$\pi_\theta(x) = \mathbf 1\{\hat\pi > \theta\}$ with
$\theta = (U_{ng} - U_{yg}) / (U_{yb} - U_{nb} + U_{ng} - U_{yg})
= 1/5 = 0.2$. The asymmetry $c_u / c_o = 4$ pushes the threshold below
one half. The retailer stocks the extra unit unless the forecast of a
stockout is very low.

\citet{GrangerPesaran2000} point out the consequence for forecast
evaluation. The realised economic value of the forecast $\hat\pi$ across
$T$ periods is
\begin{equation}
V(\hat\pi)
= T^{-1} \sum_{t=1}^T H \bigl(y_t - \theta\bigr)
  \mathbf 1\{\hat\pi_t > \theta\}
  + \text{constant},
\label{eq:gp_value}
\end{equation}
with $H = U_{yb} - U_{nb} + U_{ng} - U_{yg} = 5 > 0$ a positive scale and
the constant invariant across forecasts. Realised value depends on
$\hat\pi$ only through the threshold-crossing indicator. Two forecasts that
look the same on mean squared error can crossover at different rates and
therefore deliver different realised values, and the statistical ranking
can reverse the decision-economic ranking.\footnote{The concrete reversal
example of \citet{GrangerPesaran2000} constructs two forecasters with the
same MSE such that one is always above $\theta$ and the other always
below. The MSEs are identical because the errors $\hat\pi^{(1)} - \pi$ and
$\hat\pi^{(2)} - \pi$ are equal in absolute value at every $t$; the
realised values differ by the full gap $H \mathbb E[|y - \theta|]$ per
period.} A forecast trained to minimise squared error is good at the wrong
thing. A forecast trained to maximise~\eqref{eq:gp_value} is good at the
right thing but typically biased on MSE.

The statistical complement of the argument is the proper scoring rule
machinery of \citet{GneitingRaftery2007}. A scoring rule $S(P, y)$
on a class of probability distributions $\mathcal P$ is proper when
$\mathbb E_{y \sim Q}[S(Q, y)] \ge \mathbb E_{y \sim Q}[S(P, y)]$
for every $P, Q \in \mathcal P$, and strictly proper when the inequality
is strict for $P \ne Q$. The Brier score \citep{Brier1950}
$S^{\mathrm B}(\hat\pi, y) = -(y - \hat\pi)^2$ is the canonical example for
binary outcomes. The \citet{GneitingRaftery2007} characterisation theorem
states that every regular proper scoring rule arises from a convex
generating function $G$ on $\mathcal P$ via
\begin{equation}
S(P, \omega)
= G(P) + G^\star(P, \omega) - \int G^\star(P, \cdot)\, \mathrm dP,
\label{eq:proper_score_rep}
\end{equation}
where $G^\star(P, \cdot)$ is a subtangent of $G$ at $P$ and the divergence
$d(P, Q) = S(Q, Q) - S(P, Q)$ is the Bregman divergence under $G$. The
Brier score corresponds to $G(p) = \|p\|_2^2$ on the simplex, the
logarithmic score to negative Shannon entropy, and the continuous ranked
probability score and the energy score to analogous constructions on richer
sample spaces.\footnote{Strict propriety prevents the forecaster from
hedging. Any rule that admits a $P \ne Q$ with the same expected score
under $Q$ gives the forecaster an incentive to misreport.}

Proper scoring rules elicit the supreme solution $\hat\pi = \pi$, the true
conditional probability, irrespective of the downstream decision threshold
$\theta$. This is convenient when the forecaster does not know the
decision maker's cost structure. It is also the source of the
Granger-Pesaran problem. The supreme forecast is the right target only
when the forecaster's hypothesis class can recover $\pi$ exactly. Under
misspecification, namely when no $\theta \in \Theta$ achieves $f_\theta(x) =
\pi(x)$, the best-in-class forecast under a strictly proper rule is the
$L^2$ projection of $\pi$ onto the hypothesis class, which is generally
not the value-maximising forecast under~\eqref{eq:gp_value}. A sequential
analogue of the same calibration question recurs in the model-based
reinforcement-learning literature, where a value-targeted surrogate loss
that is itself uncalibrated can return the wrong transition model and the
wrong value function even at the population level
\citep{VoelckerCalib2025}, a point we take up in
Section~\ref{sec:vaml_ve}.

\citet{Dawid1998} closes the loop between decision theory and proper scoring
rules. Given an action space $\mathcal A$, a utility $U(\omega, a)$, and the
Bayes act $a_P = \arg\max_a \mathbb E_P[U(\omega, a)]$ under forecast
distribution $P$, the score $S(P, \omega) = U(\omega, a_P)$ is proper with
respect to the convex class containing $P$, and strictly proper when the
Bayes act is unique. The decision-induced score depends on the entire utility
structure of the downstream problem, while a generic strictly proper rule
such as CRPS does not. The chapter's working definition of decision-focused
learning is the construction that uses Dawid's reduction in reverse. Pick the
utility $U$ first, derive the induced score, and train the forecaster against
that score. Wald (1950) statistical decision theory carries the prescription
in the abstract. Sections~\ref{sec:dfl}-\ref{sec:adp} carry it out in
practice for single-period and sequential settings respectively.

What loss should we train under? The supreme answer is a strictly proper
rule on the conditional distribution. The pragmatic answer is the
decision-induced score $S(P, \omega) = U(\omega, a_P)$ when the downstream
decision and utility are known. The next section takes up the pragmatic
case in the form that scales to high-dimensional predictors.
